
\begin{ledgroupsized}[r]{120mm}
\footnotesize
\pstart
\noindent\textbf{Überlieferung:}
\pend
\end{ledgroupsized}
\begin{ledgroupsized}[r]{114mm}
\footnotesize
\pstart
\parindent -6mm
\makebox[6mm][l]{\textit{L}}%
Auszüge mit Bemerkungen aus
\indexN{\textls{Varignon}, Pierre de 1654\textendash1722}\textsc{P.~Varignon},
\cite{02028}\textit{Projet d'une nouvelle mechanique avec un examen de l'opinion de M.\;Borelli, sur les proprietéz des poids suspendus par des cordes}, Paris 1687:
LH~XXXVII~5 Bl.~217\textendash218. 
Ein Bogen~2\textsuperscript{o};
\href{\Scan{LH_37_05_218wz}}%
{Wasserzeichen} in Bl.~218 (italienisches Papier);
Papiererhaltungsmaßnahmen;
alle Ränder beschnitten; geringfügiger Textverlust an den Bogenaußenrändern durch Papierschaden sowie infolge der Papiererhaltungsmaßnahmen
(einige der verlorenen Textteile sind in dem 1967 angefertigten S-Film Nr.~116 der GWLB Hannover noch lesbar).
Vier Seiten.
\pend
\end{ledgroupsized}
\vspace{5mm}
\begin{ledgroup}
\footnotesize
\pstart
\noindent
\textbf{Datierungsgründe:}
Mit N.\;\ref{60365} liegen Leibnizens Auszüge aus zwei \cite{02028}Werken über die Statik vor, die 
\indexN{\textls{Varignon}, Pierre de 1654\textendash1722}Pierre Varignon 
um 1687 veröffentlichte. 
Seiner Arbeit über den Grundsatz des Parallelogramms der Kräfte (\textit{Projet d'une nouvelle mechanique})
war eine Kritik eines statischen Theorems in \indexN{\textls{Borelli} (Borellus), Giovanni Alfonso 1608\textendash1679}G.\,A.~Borellis \cite{02099}\textit{De motu animalium}, Rom 1680 beigebunden (\textit{Examen}).
Leibniz hat aus beiden Teilen von \indexN{\textls{Varignon}, Pierre de 1654\textendash1722}Varignons
\cite{02028}Publikation exzerpiert, wobei der Großteil seiner Auszüge (\refpassage{37_05_217-218_4a}{37_05_217-218_4b}) auf das \textit{Projet} (S.\;1\textendash83 in der Vorlage) entfällt, 
während nur das untere Drittel der letzten Seite des Bogens (Bl.~218~v\textsuperscript{o}, hier \refpassage{37_05_217-218_5a}{37_05_217-218_5b}) dem \textit{Examen} (S.\;85\textendash133 bei \indexN{\textls{Varignon}, Pierre de 1654\textendash1722}Varignon) gewidmet ist,
und zwar lediglich dem Abschnitt auf S.\;92\textendash100.
\pend
\pstart
Die vorgeschlagene Datierung beruht auf dem Papier (italienischer Herstellung) und dem Wasserzeichen (Wappen der Medici), das bei Leibniz 
für die Zeit seiner Aufenthalte in \indexO{Florenz}Florenz (1. \textendash\ 22. Dezember 1689) 
und Modena\indexO{Modena} (30. Dezember 1689 \textendash\ 2. Februar 1690) belegt ist.
Die kumulative Zeitspanne dieser Aufenthalte bildet den wahrscheinlichen Entstehungszeitraum der Auszüge.
Mangels ausdrücklicher Belege einer Auseinandersetzung mit \indexN{\textls{Varignon}, Pierre de 1654\textendash1722}Varignons Publikation in der Korrespondenz der italienischen Zeit erscheint es nicht möglich, die Datierung näher zu präzisieren.
\pend
\end{ledgroup}
\newpage
\selectlanguage{french}
\vspace{8mm}
\pstart
\centering
\edlabel{37_05_217-218_4a}%
\href{\Scan{LH_37_05_217r}}%
{\lbrack217~r\textsuperscript{o}\rbrack}\ 
\hfill
\edtext{%
\textit{Projet d'une nouvelle mechanique avec un} 
\edtext{\textit{examen de}}{%
\lemma{}%
\Bfootnote{%
\textit{examen} \bars{\textit{de l'opinion} \textit{streicht Hrsg.}} \textit{de}~\textit{L}}}
\textit{l'opinion de} 
\hfill
\phantom{\lbrack217~r\textsuperscript{o}\rbrack}
\linebreak
\textit{M.\;Borelli\indexN{\textls{Borelli} (Borellus), Giovanni Alfonso 1608\textendash1679} sur les proprietés des poids suspendus par des cordes.}%
}{%
\lemma{\textit{Projet} \lbrack...\rbrack\ \textit{cordes}}%
\Cfootnote{
\indexN{\textls{Varignon}, Pierre de 1654\textendash1722}\textsc{P.~Varignon}, \cite{02028}\textit{Projet d'une nouvelle mechanique avec un examen de l'opinion de M.\;Borelli, sur les proprietéz des poids suspendus par des cordes}, Paris 1687, Titelseite.}}
Paris 1687
\pend
\vspace{0.5\baselineskip}
\pstart
\noindent
\edtext{Preface.}{\lemma{}\Bfootnote{Preface. \textit{erg.~L}}}
\edtext{Pour \textit{chercher l'equilibre\indexS{equilibre} dans sa generation}}{%
\lemma{Pour \lbrack...\rbrack\ \textit{generation}}%
\Cfootnote{\cite{02028}a.\,a.\,O., Preface, S.~\lbrack IX\rbrack\ mit Änderungen.}}
\edtext{%
je consideray \textit{un poids qu'une puissance\indexS{puissance} soutient sur un plan incliné,\indexS{plan incliné} 
d'abord je me le representay de telle figure que le concours de sa ligne de direction\indexS{ligne de direction}\indexS{concours des lignes de direction}
avec celle de cette puissance\indexS{puissance} se fit dans} 
\edtext{\textit{quel}\lbrack\textit{qu'}\rbrack\textit{un}}{%
\lemma{}%
\Bfootnote{quelcun \textit{L ändert Hrsg. nach Vorlage}}}
\textit{de ses points. De là je vis que leur concours d'action se faisant aussi par} 
le \textit{moyen} dans \textit{ce seul point, il devenoit alors son centre de direction,\indexS{centre de direction}
de sorte que si ce plan eut manqué tout d'un coup ce corps auroit necessairement suivi
l'impression de ce point. Je cherchai ensuite quelle devoit estre cette impression, 
et j'apperçus que celles qui} se
\textit{faisoient sur ce point, et la pesanteur de ce poids,
et la puissance\indexS{puissance} qui le retenoit estant les memes que s'il eut esté poussé 
en meme temps par deux forces qui leur eussent esté egales et qui eussent agi
suivant leurs lignes de direction,\indexS{ligne de direction} j'apperçus dis-je qu'il luy en resultoit une impression composée,\indexS{impression composée}
suivant une ligne, qui estoit la diagonale d'un parallelogramme\indexS{diagonale du parallelogramme} fait sous des parties
de ces lignes de direction,\indexS{ligne de direction} qui estoient entre elles comme ce poids et cette puissance\indexS{puissance}}\,\lbrack\!;\rbrack\
\textit{d'où je vis que l'impression de ce corps se faisoit alors suivant cette diagonale,
qui devenoit en ce cas sa ligne de direction.\indexS{ligne de direction} Mais que ce plan 
luy estant perpendiculairement opposé il la soutenoit toute entiere,
ce qui faisoit que ce poids ainsi poussé par le concours d'action de sa pesanteur
et de la puissance,\indexS{puissance} qui luy estoit appliquée}\lbrack,\rbrack\
\textit{demeurant sur ce plan incliné\indexS{plan incliné} de meme que s'il eust esté horizontal,\indexS{plan horizontal}
et que cette impression composée\indexS{impression composée} n'eut esté qu'un effet de sa seule pesanteur.}
\textit{De cette pensée j'en vis naistre plusieurs autres, et je m'apperçus}\,\lbrack\!\!:\rbrack\
\textit{1\textsuperscript{o}\! que toute l'impression que ce plan recevoit alors de ce poids}
et \textit{cette puissance\indexS{puissance} se faisoit suivant cette diagonale.\indexS{diagonale du parallelogramme}
2\textsuperscript{o}\! que sa charge}\lbrack,\rbrack\ \textit{c'estoit à dire la force
de cette meme impression}\lbrack,\rbrack\ \textit{estoit à ce poids et à cette puissance\indexS{puissance}
comme cette meme diagonale\indexS{diagonale du parallelogramme} à chacun des costés qui les representent
dans son parallelogramme.
3\textsuperscript{o}\! que ce poids et cette puissance\indexS{puissance} estoient tousjours entre eux 
comme ces memes costés}\lbrack,\rbrack\ \textit{c'est à dire en raison reciproque
des sinus des angles\indexS{sinus de l'angle} que font leur lignes de direction avec cette diagonale,\indexS{diagonale du parallelogramme} ou
(ce qui revient au meme) en raison reciproque des distances de quelque point 
que ce soit de cette diagonale à leur lignes de direction.}}{%
\lemma{je consideray \lbrack...\rbrack\ \textit{direction}}%
\Cfootnote{\cite{02028}a.\,a.\,O., S.~\lbrack IX\textendash XI\rbrack\ mit Änderungen.}}
\pend
\pstart
\edtext{\textit{Je cherchay par le meme chemin comment des poids soutenus avec des cordes\indexS{poids soutenus avec des cordes}\indexS{corde}
ou appliqués à des poulies\indexS{poulie} ou bien à des leviers\indexS{levier} font equilibre\indexS{equilibre} entre eux 
ou avec les puissances\indexS{puissance} qui les soustiennent, et j'apperçus de meme 
que tout cela se faisoit encor par la voye des mouvemens composés.\indexS{mouvement composé}}}{%
\lemma{\textit{Je} \lbrack...\rbrack\ \textit{composés}}%
\Cfootnote{\cite{02028}a.\,a.\,O., S.~\lbrack XII\rbrack\ mit Auslassung.}}
\edtext{\textit{Ensuite au lieu} d'une seule \textit{demonstration qu'on ne fait ordinairement 
que pour les plans inclinés,\indexS{plan incliné} j'en trouvay une qui s'étend generalement à toute sorte de surfaces
et à toutes les directions possibles des puissances,\indexS{puissance} ou des poids qui y sont appliqués.}}{%
\lemma{\textit{Ensuite} \lbrack...\rbrack\ \textit{appliqués}}%
\Cfootnote{\cite{02028}a.\,a.\,O., S.~\lbrack XIII\rbrack\ mit Änderungen.}}
\edtext{Cette \textit{route facilite extremement le calcul}\lbrack,\rbrack\
les rapports des forces\indexS{rapport des forces} estant 
\textit{determinés immediatement par les sinus des angles,\indexS{sinus de l'angle} que font leur lignes de direction\indexS{ligne de direction}}
\edtext{\textit{avec celle}}{%
\lemma{\textit{avec}}%
\Bfootnote{%
\textit{(1)}~celles
\textit{(2)}~\textit{celle}~\textit{L}}}
\textit{de l'impression qui resulte de leur concours d'action, et que cette methode determine
pour le point ou elles concourent. On y voit que lors que deux puissances\indexS{puissance} 
ou deux poids ou bien une puissance\indexS{puissance} et un poids font equilibre\indexS{equilibre} soit avec des cordes\indexS{corde} seulement,
soit à l'aide de quelque poulie,\indexS{poulie} de quelque surface}
\edtext{\textit{ou}}{\lemma{}\Bfootnote{\textit{ou} \textit{erg.~L}}}
\textit{de quelque levier\indexS{levier} que ce soit, ils sont toujours entre eux 
en raison reciproque des sinus\indexS{sinus de l'angle} de ces memes angles.}}{%
\lemma{Cette \lbrack...\rbrack\ \textit{angles}}%
\Cfootnote{\cite{02028}a.\,a.\,O., S.~\lbrack XIV\rbrack\ mit Änderungen.}}
\pend
\pstart
\edtext{L'auteur\indexN{\textls{Varignon}, Pierre de 1654\textendash1722} 
espere \textit{d'expliquer} par là \textit{les effects les plus surprenans
des Machines composées\indexS{Machine composée} de l'art et de la nature.}}{%
\lemma{L'auteur \lbrack...\rbrack\ \textit{nature}}%
\Cfootnote{\cite{02028}a.\,a.\,O., S.~\lbrack XIV\rbrack\ mit Änderungen.}}
\pend
\pstart
\edtext{\textit{\textls{Axiome}. Les espaces que parcourent des corps egaux 
dans des temps egaux sont comme les forces\indexS{force} et reciproquement.}}{%
\lemma{\textit{\textls{Axiome}.} \lbrack...\rbrack\ \textit{reciproquement}}%
\Cfootnote{\cite{02028}a.\,a.\,O., S.\;1 mit Auslassungen.}}
\pend
\pstart
\edtext{\textit{\textls{Demande}. On suppose que dans tout corps qui se meut 
ou qui fait effort pour se mouvoir il y a tousjours un certain point qui surchargé
de l'impression de tous les autres determine ce corps à suivre celle qu'il a pour lors
vers l'endroit où il tend.
\textls{On ne se met point en peine} que ce point soit
\textls{le meme dans toutes les situations possibles de ce corps}.
C'est assez que dans chaque situation il y en ait un, qu'on appelle icy
son \textls{centre de gravité\indexS{centre de gravité}} or plus generalement 
son centre de direction\indexS{centre de direction} ou d'equilibre,\indexS{centre d'equilibre}\indexS{equilibre} du moins pour le temps 
qu'il determine ainsi ce corps à suivre son impression, et la ligne qui joint ce point
avec celuy où il tend, s'appelle la ligne de direction.\indexS{ligne de direction}}}{%
\lemma{\textit{\textls{Demande}.} \lbrack...\rbrack\ \textit{direction}}%
\Cfootnote{\cite{02028}a.\,a.\,O., S.\;2 mit Änderungen.}}
\pend
\begin{edfigure}%
\hspace{1em}%
\includegraphics[scale=1]{%
gesamttex/edit_VIII-4/images/LH_37_05_217-218_d1_217r.pdf}
\hspace*{15mm}%
\includegraphics[scale=1]{%
gesamttex/edit_VIII-4/images/LH_37_05_217-218_d2_217r.pdf}
\edcaption{\edtext{\lbrack\textit{Fig.~1}\rbrack}{%
\lemma{\lbrack\textit{Fig.~1}\rbrack}%
\Cfootnote[nonum]{\cite{02028}a.\,a.\,O., pl.~1, fig.~3; die zum \textit{Projet} gehörigen Tafeln befinden sich auf neun Durchschussblättern am Ende dieses Werks, nach S.\;83.}}%
\hspace*{25mm}%
\edtext{\lbrack\textit{Fig.~2}\rbrack}{%
\lemma{\lbrack\textit{Fig.~2}\rbrack}%
\Cfootnote[nonum]{\cite{02028}a.\,a.\,O., pl.~1, fig.~4. Die Strecke \textit{NT} fehlt in der Vorlage.}%
\lemma{}%
\Afootnote[nonum]{\selectlanguage{latin}%
\textit{Neben dem Diagramm \lbrack\textit{Fig.~2}\rbrack}:
Sed quid si palus \textit{HN} continuatus
sustineatur in duobus locis ut in \textit{T} et \textit{H} ut non possit descendere \textit{A} per \textit{AB}.
\textit{HM} et \textit{HN} pali ut ab autore\indexN{\textls{Varignon}, Pierre de 1654\textendash1722}
exhibentur inutiles, nam in nullo situ sustinent.}}}
\end{edfigure}
\pstart
\edtext{%
\edtext{\textit{Lemme~1.}}{\lemma{}\Bfootnote{\textit{Lemme~1.} \textit{erg.~L}}}
\textit{Le poids\indexS{poids suspendu de deux cordes} \textit{A} suspendu de deux cordes\indexS{corde} \textit{MH}, \textit{NH}, ou soutenu de deux} 
\edtext{\textit{pieux} se mouvera \textit{jusqu'à ce que}}{%
\lemma{\textit{pieux}}%
\Bfootnote{%
\textit{(1)}~sera en repos, si
\textit{(2)}~se mouvera \textit{jusqu'à ce que}~\textit{L}}}
\textit{la} droite \textit{\textit{AH} qui joint le point de suspension\indexS{point de suspension}} \textit{H}, 
\edtext{\textit{avec le centre}}{%
\lemma{avec le}%
\Bfootnote{%
\textit{(1)}~point de dir
\textit{(2)}~centre~\textit{L}}} 
\textit{de} direction\indexS{centre de direction} \textit{\textit{A} soit dans la perpendiculaire,} ou pour les 
\edtext{\lbrack p\rbrack ieux}{\lemma{}\Bfootnote{epieux \textit{L ändert Hrsg. nach Vorlage}}}
dans l'horizontale, et alors il sera en repos.}{%
\lemma{\textit{Lemme~1.} \lbrack...\rbrack\ repos}%
\Cfootnote{\cite{02028}a.\,a.\,O., S.\;3 mit Änderungen.}}
(+~Les 
\edtext{\lbrack p\rbrack ieux}{\lemma{}\Bfootnote{epieux \textit{L ändert Hrsg. nach Vorlage}}}
ne repondent pas assez, car le corps 
\edtext{\textit{A}}{\lemma{}\Bfootnote{\textit{A} \textit{erg.~L}}}
arrive en (\textit{B}) peut aller encor plus bas.~+)
\pend
\begin{edfigure}[1em]%
\includegraphics[scale=1]{%
gesamttex/edit_VIII-4/images/LH_37_05_217-218_d3_217r.pdf}
\edcaption{\edtext{\lbrack\textit{Fig.~3}\rbrack}{%
\lemma{\lbrack\textit{Fig.~3}\rbrack}%
\Cfootnote[nonum]{\cite{02028}a.\,a.\,O., pl.~1, fig.~5.}}}
\end{edfigure}
\pstart
\edtext{%
\textit{Lemm.~2. Le poids\indexS{poids suspendu de deux cordes} \textit{MN} estant suspendu par 2 cordes\indexS{corde} \textit{PM} et \textit{RN}}
\edtext{\textit{attach}\lbrack\textit{é}\rbrack\textit{es}}{%
\lemma{}\Bfootnote{attaches \textit{L ändert Hrsg. nach Vorlage}}}
\textit{aux clous \textit{P} et \textit{R}, et qui prolongées concourent en \textit{H},
sa ligne de direction\indexS{ligne de direction} \textit{AH} passera par ce point de concours.}
Nam idem est 
\edtext{\lbrack utrum\rbrack}{\lemma{}\Bfootnote{utrum \textit{erg.~Hrsg.}}}
chorda\indexS{chorda} sit \textit{MP} an \textit{MH}, vel \textit{NR} aut \textit{NH} jam 
\edtext{si chordae\indexS{chorda}}{%
\lemma{si}%
\Bfootnote{%
\textit{(1)}~sint
\textit{(2)}~chordae~\textit{L}}}
concurrant in \textit{H}, linea 
\edtext{directio\lbrack nis\rbrack\indexS{linea directionis}}{%
\lemma{}%
\Bfootnote{directio \textit{L ändert Hrsg.}}}
est \textit{AH} per lemma~1. Ergo et nunc.}{%
\lemma{\textit{Lemm.~2.} \lbrack...\rbrack\ nunc}%
\Cfootnote{\cite{02028}a.\,a.\,O., S.\;4\,f. mit Änderungen.}}
\pend
\pstart
\edtext{\textit{Lemm.~3. Si un point \textit{A} sans pesanteur est poussé en meme temps 
par deux puissances\indexS{puissance} \textit{E} et \textit{F} suivant les lignes}
\edtext{\textit{AC}.\lbrack\textit{AB}\rbrack}{%
\lemma{}%
\Bfootnote{%
\textit{AC}.\textit{BC} 
\textit{L ändert Hrsg. nach Vorlage}}}
\textit{qui fassent un angle \textit{CAB}, et}
\edtext{\textit{soit puissance\indexS{puissance}}}{%
\lemma{\textit{soit}}%
\Bfootnote{%
\textit{(1)}~\textit{E} à \textit{F} \bars{vel \textit{streicht Hrsg.}} ut
\textit{(2)}~\textit{puissance}~\textit{L}}}
\textit{\textit{E} à puissance\indexS{puissance} \textit{F}
comme \textit{AC} à \textit{AB}, ce point \textit{A} suivra la diagonale\indexS{diagonale du parallelogramme} \textit{AD}
du parallelogramme \textit{BC} fait sous ces deux lignes.}}{%
\lemma{\textit{Lemm.~3.} \lbrack...\rbrack\ lignes}%
\Cfootnote{\cite{02028}a.\,a.\,O., S.\;6 mit Änderungen.}}
\pend
\pstart
\edtext{Lemme~4. eodem
\edtext{re\lbrack dd\rbrack it.}{\lemma{}\Bfootnote{redidit \textit{L ändert Hrsg.}}}}{%
\lemma{Lemme~4. eodem re\lbrack dd\rbrack it}%
\Cfootnote{Siehe \cite{02028}a.\,a.\,O., S.\;7\,f.}}
\pend
\pstart
\edtext{\textit{Lemme~5. Les trois costés d'un triangle rectiligne quelqu'il soit sont entre eux 
comme les sinus des angles opposés.\indexS{sinus de l'angle opposé}}}{%
\lemma{\textit{Lemme~5.} \lbrack...\rbrack\ \textit{opposés}}%
\Cfootnote{\cite{02028}a.\,a.\,O., S.\;9 mit Auslassung.}}
\pend
\begin{edfigure}[0.5em][1em]
\hspace{1.5em}%
\includegraphics[scale=1]{%
gesamttex/edit_VIII-4/images/LH_37_05_217-218_d4_217r.pdf}
\edcaption{\edtext{\lbrack\textit{Fig.~4}\rbrack}{%
\lemma{\lbrack\textit{Fig.~4}\rbrack}%
\Cfootnote[nonum]{%
\cite{02028}a.\,a.\,O., pl.~2, fig.~14. Ein gestrichener Entwurf zum Diagramm wird nicht wiedergegeben.}}}
\end{edfigure}
\pstart
\edtext{%
\textit{\textls{Proposition fondamentale des poids suspendus avec des cordes\indexS{poids suspendu avec des cordes}}.}
Poids ou puissance\indexS{puissance} \textit{K} est à la puissance\indexS{puissance} \textit{R} (vel \textit{P})
qui soutient la corde\indexS{corde} \textit{RA} (\textit{PA}) comme le sinus de l'angle\indexS{sinus de l'angle} \textit{PAR}
au sinus de l'angle\indexS{sinus de l'angle} \textit{PAD}
\edtext{(\textit{RAD}).
\selectlanguage{latin}%
Et}{%
\lemma{(\textit{RAD})}%
\Bfootnote{%
\textit{(1)}~.~scilicet
\textit{(2)}~Et~\textit{L}}}
potentiae\indexS{potentia} proinde \textit{R} et \textit{P} sunt reciproce ut sinus angulorum\indexS{sinus anguli} 
quos directiones earum faciunt in directione potentiae sustentatae 
seu quae cum ipsis in aequilibrio\indexS{aequilibrium} est. Idem enim est ac si 
\edtext{punctum \textit{A} duabus potentiis\indexS{potentia}}{%
\lemma{punctum \textit{A}}%
\Bfootnote{%
\textit{(1)}~duobus
\textit{(a)}~nisibus impelleretur
\textit{(b)}~conatibus impe
\textit{(2)}~duabus potentiis~\textit{L}}}
(+~ut ego loquor conatibus\indexS{conatus}~\lbrack+\rbrack)
impelleretur \textit{AC}, \textit{AB}, sunt autem \textit{AC}\lbrack,\rbrack\ \textit{AB}, \textit{AD} 
ut sinus angulorum oppositorum\indexS{sinus anguli oppositi} id est dicto modo.}{%
\lemma{\textls{\textit{Proposition}} \lbrack...\rbrack\ modo}%
\Cfootnote{\cite{02028}a.\,a.\,O., S.\;11\textendash13 mit Änderungen.}}
\edtext{\selectlanguage{french}%
Comme les sinus sont comme les costés du triangle \textit{ADC} vel \textit{ABD},
et les puissances\indexS{puissance} de meme, tousjours deux de ces puissances\indexS{puissance} 
seront plus grandes que la 3\textsuperscript{me}.}{%
\lemma{\textls{Comme} \lbrack...\rbrack\ 3\textsuperscript{me}}%
\Cfootnote{\cite{02028}a.\,a.\,O., siehe Coroll.~III, S.\;14.}}
\pend
\pstart
\edlabel{37_05_217-218_7a}%
\edtext{}{%
\xxref{37_05_217-218_7a}{37_05_217-218_7b}%
\lemma{\textls{S'il} \lbrack...\rbrack\ concourent}%
\Cfootnote{Siehe \cite{02028}a.\,a.\,O., Coroll.~VII, S.\;16\,f.}}%
\textls{S'il y avoit} plusieurs cordes\indexS{corde} on 
\edtext{peut prendre pour une seule}{%
\lemma{peut}%
\Bfootnote{%
\textit{(1)}~conclure
\textit{(2)}~prendre pour une seule~\textit{L}}}
l'impression
\edtext{composée\indexS{impression composée}}{\lemma{}\Bfootnote{composée \textit{erg.~L}}}
de chacun des points où deux de ces cordes\indexS{corde} 
\edtext{concourent\edlabel{37_05_217-218_7b}
\selectlanguage{latin}%
(+~debet}{%
\lemma{concourent (+}%
\Bfootnote{%
\textit{(1)}~potest
\textit{(2)}~debet~\textit{L}}}
autem concipi pondus totum esse redactum in punctum \textit{A} ubi chordae\indexS{chorda} concurrunt. 
Hoc autem punctum concursus semper esse in linea directionis.\indexS{linea directionis}~\lbrack+)\rbrack\
Autor demonstrare
\edtext{voluit lemm.~1.}{%
\lemma{voluit}%
\Bfootnote{%
\textit{(1)}~prop.~1
\textit{(2)}~lemm.~1.~\textit{L}}}
Sed demonstratio ejus non applicatur ad casum, 
\edtext{quo intersectio chordarum\indexS{chorda}\indexS{intersectio chordarum} est infra.}{%
\lemma{quo}%
\Bfootnote{%
\textit{(1)}~corpora
\textit{(2)}~intersectio chordarum
\textit{(a)}~\bars{est \textit{streicht Hrsg.}} infra pondus
\textit{(b)}~est infra.~\textit{L}}}
Videtur tamen et sic demonstrari posse fingendo chordas\indexS{chorda} \textit{R} et \textit{P}
revera concurrere in \textit{A}, et 
\edtext{\lbrack\textit{K}\rbrack}{\lemma{}\Bfootnote{\textit{A} \textit{L ändert Hrsg.}}}
suspensum esse
\edtext{in puncto}{%
\lemma{in}%
\Bfootnote{%
\textit{(1)}~linea
\textit{(2)}~puncto~\textit{L}}}
lineae directionis\indexS{linea directionis}
\edtext{infra \textit{A}, ut}{%
\lemma{infra \textit{A}}%
\Bfootnote{%
\textit{(1)}~ex
\textit{(2)}~ut~\textit{L}}}
in 
\edtext{\lbrack\textit{H}\rbrack,}{\lemma{}\Bfootnote{\textit{K} \textit{L ändert Hrsg.}}}
utique chorda\indexS{chorda} \textit{H} annectens chordae\indexS{chorda} \textit{RAP}, annectetur in \textit{A}.
Est autem \textit{AH} in linea directionis.\indexS{linea directionis} Ergo et nunc \textit{A} est in linea directionis.\indexS{linea directionis}
\pend
\begin{edfigure}[0.75em][0.5em]
\includegraphics[scale=1]{%
gesamttex/edit_VIII-4/images/LH_37_05_217-218_d5_217r.pdf}
\edcaption{\edtext{\lbrack\textit{Fig.~5}\rbrack}{%
\lemma{\lbrack\textit{Fig.~5}\rbrack}%
\Cfootnote[nonum]{%
Ein gestrichener Entwurf zum Diagramm wird nicht wiedergegeben.}
\lemma{}\Afootnote[nonum]{\selectlanguage{latin}%
\textit{Unterhalb der Zeichnung \lbrack\textit{Fig.~5}\rbrack}: 
Quid si duo sint pali sustinentes\textsuperscript{\lbrack a\rbrack} 
corpus \textit{FMA}, nempe \textit{MP}, et \textit{RN}, sed ita ut concursus eorum \textit{H}
non fia\lbrack t\rbrack\textsuperscript{\lbrack b\rbrack} in linea directionis\indexS{linea directionis} \textit{FC} per centrum gravitatis\indexS{centrum gravitatis} transiente.
Oportet autem palos esse perpendiculares ad superficiem corporis, alioqui non sustinerent.
\\[1.5ex] %
{\footnotesize 
\textsuperscript{\lbrack a\rbrack}~sustinentes \textit{(1)}~centrum \textit{(2)}~corpus~\textit{L}\quad 
\textsuperscript{\lbrack b\rbrack}~fiant \textit{L ändert Hrsg.}}}}%
}
\end{edfigure}
\pstart
\noindent
\href{\Scan{LH_37_05_217v}}%
{\lbrack217~v\textsuperscript{o}\rbrack}
\pend
\nobreak
\begin{edfigure}[0.0em][1.0em]
\includegraphics[scale=1]{%
gesamttex/edit_VIII-4/images/LH_37_05_217-218_d6_217v.pdf}
\edcaption{\edtext{\lbrack\textit{Fig.~6}\rbrack}{%
\lemma{\lbrack\textit{Fig.~6}\rbrack}%
\Cfootnote[nonum]{\cite{02028}a.\,a.\,O., pl.~3, fig.~20. Die von Leibniz gegenüber der Vorlage vertauschten Bezeichnungen der Punkte \textit{P} und \textit{R} werden vom Hrsg. wiederhergestellt.}}}
\end{edfigure}
\selectlanguage{french}
\pstart
\edtext{%
\textit{Proposition fondamentale \textls{des poulies}.\indexS{poulie}
Soit la puissance\indexS{puissance} ou le poids \textit{D} appliqué ou suspendu au centre mobile\indexS{centre mobile de la poulie}
d'une poulie \textit{A} autour de la quelle passe la corde\indexS{corde} \textit{PMNR}, dont les extremités
sont retenues par les puissances\indexS{puissance} \textit{P} et \textit{R}: quelque angle \textit{MHN} 
que fassent entre elles les parties \textit{MP} et \textit{NR} de cette corde,\indexS{corde} 
prolongées jusqu'à ce qu'elles concourent en \textit{H}}\,\lbrack\!;\rbrack\ 
\textit{le poids ou la puissance\indexS{puissance} \textit{D} sera tousjours à chacune des puissances\indexS{puissance} \textit{P} et \textit{R},
comme le sinus de cet angle\indexS{sinus de l'angle} au sinus de la}
\edtext{\textit{moitié.} Car \textit{les puissances\indexS{puissance} \textit{P} et \textit{R}} estant \textit{en equilibre\indexS{equilibre}} entre}{%
\lemma{\textit{moitié.}}%
\Bfootnote{%
\textit{(1)}~\textit{Les puiss}
\textit{(2)}~\textit{Car les puissances}
\textit{(a)}~\textit{estant}
\textit{(b)}~\textit{\textit{P} et \textit{R} estant en equilibre}
\textit{(aa)}~aussi bien
\textit{(bb)}~entre~\textit{L}}}
elles et \textit{avec la puissance\indexS{puissance} \textit{D}, 
la corde\indexS{corde} \textit{PNMR} demeure fixe dessus sans glisser,}
et c'est autant \textit{que si elle y estoit collée}.}{%
\lemma{\textit{Proposition} \lbrack...\rbrack\ \textit{collée}}%
\Cfootnote{\cite{02028}a.\,a.\,O., S.\;25 mit Änderungen.}}
\edtext{%
C'est donc la meme chose \textit{que si \textit{PN} et \textit{MR} estoient deux cordes\indexS{corde}} differentes \textit{attachées} 
\edtext{\textit{separément} à}{%
\lemma{\textit{separément}}%
\Bfootnote{%
\textit{(1)}~au \textit{NM}
\textit{(2)}~à~\textit{L}}}
la poulie\indexS{poulie} \textit{NM}. \textit{Or en ce cas la ligne de direction\indexS{ligne de direction}}
\edtext{\textit{AD}}{\lemma{}\Bfootnote{\textit{AD} \textit{erg.~L}}}
de la puissance\indexS{puissance} \textit{D} \textit{passera par \textit{H}, point} du concours
\textit{(Lemm.~2 cor.~1}\lbrack\textit{)}\rbrack\
\textit{mais encor cette poulie\indexS{poulie} regardée avec une telle impression, seroit
(\textup{par la} prop. fondamentale des} 
\edtext{\textit{cordes\indexS{corde}) aux}}{%
\lemma{\textit{cordes)}}%
\Bfootnote{%
\textit{(1)}~comme
\textit{(2)}~\textit{aux}~\textit{L}}}
\textit{puissances\indexS{puissance} \textit{P} et \textit{R} qui la retiennent comme le sinus de l'angle\indexS{sinus de l'angle} \textit{MHN}}
\edtext{\textit{au}\lbrack\textit{x}\rbrack}{%
\lemma{}\Bfootnote{\textit{au}\lbrack\textit{x}\rbrack\ \textit{L ändert Hrsg. nach Vorlage}}}
\textit{sinus des angles\indexS{sinus de l'angle} \textit{NHA}, et}
\edtext{\textit{MHA} qui}{%
\lemma{\textit{MHA}}%
\Bfootnote{%
\textit{(1)}~la m
\textit{(2)}~or les mo
\textit{(3)}~qui~\textit{L}}}
\textit{sont icy les moitiés de l'angle \textit{MHN}.}}{%
\lemma{C'est \lbrack...\rbrack\ \textit{MHN}}%
\Cfootnote{\cite{02028}a.\,a.\,O., S.\;26 mit Änderungen.}}
\edtext{La meme chose se doit entendre de la resistence de \textit{D}, si \textit{D} estoit seulement un clou.}{%
\lemma{La meme \lbrack...\rbrack\ clou}%
\Cfootnote{\cite{02028}a.\,a.\,O., Coroll.~V, S.\;28.}}
\pend
\pstart
\edtext{La puissance\indexS{puissance} \textit{D} n'est double des puissances\indexS{puissance} \textit{P} et \textit{R}, 
que lors que \textit{PN} et \textit{MR} sont 
\edlabel{37_05_217-218_1a}%
\edtext{}{%
\xxref{37_05_217-218_1a}{37_05_217-218_1b}%
\lemma{paralleles;}%
\Bfootnote{%
\textit{(1)}~lors ce
\textit{(2)}~ce que lors ceux gr
\textit{(3)}~ce~\textit{L}}}%
paralleles;}{%
\lemma{La puissance \lbrack...\rbrack\ paralleles}%
\Cfootnote{\cite{02028}a.\,a.\,O., Coroll.~VIII, S.\;30.}}
\edtext{ce\edlabel{37_05_217-218_1b} qu'on a pourtant regardé comme general.
\textit{Il est vray que les demonstrations des auteurs} sont fondées sur le parallelisme des cordes,
mais ayant enoncé la conclusion generalement, ils en ont abusé
quelques fois sans se souvenir de la demonstration qui est limitée.}{%
\lemma{ce qu'on a \lbrack...\rbrack\ limitée}%
\Cfootnote{\cite{02028}a.\,a.\,O., S.\;31\,f. mit Änderungen.}}
\pend
\begin{edfigure}[0.5em][1em]
\includegraphics[scale=1]{%
gesamttex/edit_VIII-4/images/LH_37_05_217-218_d7_217v.pdf}
\edcaption{\edtext{\lbrack\textit{Fig.~7}\rbrack}{%
\lemma{\lbrack\textit{Fig.~7}\rbrack}%
\Cfootnote[nonum]{Teilweise Übernahme der Vorlage, \cite{02028}a.\,a.\,O., pl.~3, fig.~24.
In der Vorlage hängt der Körper \textit{D} an einer dritten Rolle \textit{A}, deren Seile einen Winkel \textit{H} bilden (siehe hier die gestrichene Variante zu \refnote{FnB:37_05_217-218_angle-H}).}}}
\end{edfigure}
\pstart
\edtext{%
S'il ya plusieurs poulies\indexS{poulie} comme \textit{C}, \textit{B},
\edtext{etc.}{\lemma{}\Bfootnote{etc. \textit{erg.~L}}}
\textit{le poids \textit{D}} est \textit{à la puissance\indexS{puissance} \textit{R}, comme le produit des sinus\indexS{sinus de l'angle} des}
\edtext{\textit{Angles \textit{K},}}{%
\lemma{\textit{Angles}}%
\applabel{FnB:37_05_217-218_angle-H}%
\Bfootnote{%
\textit{(1)}~\textit{H}
\textit{(2)}~\textit{\textit{K},}~\textit{L}}}
\textit{\textit{L}, etc. qui}
\edtext{\textit{font les cordes\indexS{corde}}}{%
\lemma{\textit{font les}}%
\Bfootnote{%
\textit{(1)}~chordes
\textit{(2)}~\textit{cordes}~\textit{L}}}
\textit{prolongées qui touchent ces poulies,\indexS{poulie}}
\textit{au produit des sinus de} 
\edtext{\textit{chacun leur}}{%
\lemma{}%
\Bfootnote{%
\textit{chacun} \bars{des \textit{streicht Hrsg.}}  \textit{leur}~\textit{L}}}
\edlabel{37_05_217-218_2a}%
\edtext{}{%
\xxref{37_05_217-218_2a}{37_05_217-218_2b}%
\lemma{\textit{moitié.}}%
\Bfootnote{%
\textit{(1)}~Les au
\textit{(2)}~Il~\textit{L}}}%
\textit{moitié.}}{%
\lemma{S'il \lbrack...\rbrack\ \textit{moitié}}%
\Cfootnote{\cite{02028}a.\,a.\,O., Coroll.~XIV, S.\;33 mit Änderungen.}}
\edtext{Il\edlabel{37_05_217-218_2b} n'y a qu'un seul cas, 
ou ce que les auteurs \textit{ont avancé} generalement,
(\textit{que le poids \textit{D} est à la puissance\indexS{puissance} \textit{R}, comme le dernier terme 
de la progression} Geometrique\indexS{progression Geometrique} \textit{double, qui en a autant qu'il y a de moufles,}
\edtext{\lbrack\textit{plus un,}\rbrack}{\lemma{}\Bfootnote{\textit{plus un,} \textit{erg.~Hrsg. nach Vorlage}}}
\textit{est au premier}\!\!:) a lieu. Sçavoir quand les parties des cordes\indexS{corde} qui touchent une meme poulie\indexS{poulie} sont paralleles.\indexS{cordes paralleles}}{%
\lemma{Il n'y a \lbrack...\rbrack\ paralleles}%
\Cfootnote{\cite{02028}a.\,a.\,O., Coroll.~XV, S.\;34 mit Änderungen.}}
\pend
\begin{edfigure}[1em][1em]
\includegraphics[scale=1]{%
gesamttex/edit_VIII-4/images/LH_37_05_217-218_d8_217v.pdf}
\edcaption{\edtext{\lbrack\textit{Fig.~8}\rbrack}{%
\lemma{\lbrack\textit{Fig.~8}\rbrack}%
\Cfootnote[nonum]{\cite{02028}a.\,a.\,O., pl.~3, fig.~26.}}}
\end{edfigure}
\pstart
\edlabel{37_05_217-218_3a}%
\edtext{}{%
\xxref{37_05_217-218_3a}{37_05_217-218_3b}%
\lemma{Que \textit{si} \lbrack...\rbrack\ \textit{dfb}}%
\Cfootnote{\cite{02028}a.\,a.\,O., Coroll.~XVI, S.\;35\,f. mit Änderungen.}}%
Que \textit{si l'on se sert de plusieurs} 
\edtext{\lbrack\textit{poulies\indexS{poulie} liées}\rbrack}{\lemma{}\Bfootnote{moufles liés \textit{L ändert Hrsg. nach Vorlage}}}
\textit{ensemble\indexS{poulies liées ensemble} comme on le voit}
\edtext{icy. \textit{La}}{%
\lemma{icy.}%
\Bfootnote{%
\textit{(1)}~Le
\textit{(2)}~\textit{La}~\textit{L}}}
\textit{puissance\indexS{puissance} \textit{R} est au poids \textit{D}, qui elle soutient à l'aide de ces poulies,\indexS{poulie}
comme \textls{le produit} des sinus des moitiés de chacun des angles que font
(si on les prolonge) les cordes\indexS{corde} tangentes des poulies mobiles,\indexS{poulie mobile} \textit{L}, \textit{K}, \textit{H}; 
à \textls{la somme des produits} de chacun des sinus\indexS{sinus de l'angle} 
des ces memes angles par le sinus des moitiés de chacun des autres.}
\textit{Soit le sinus de l'angle\indexS{sinus de l'angle} \textit{A}}\lbrack,\rbrack\ \textit{(\textit{a}), de sa moitié, \textit{b}; 
de l'angle \textit{C}, (\textit{c}) de sa moitié \textit{D}, de l'angle \textit{E}, (\textit{e}),
de sa moitié \textit{f}; je dis que la puissance\indexS{puissance} \textit{R}, est au poids \textit{D},
comme \textit{bdf}} est \textit{à $(a)df + (c)bf + (e)bd$. Parce que la corde\indexS{corde} \textit{RRORNRM}
estant egalement}
\edtext{\textit{bandé}\lbrack\textit{e}\rbrack}{\lemma{}\Bfootnote{bandées \textit{L ändert Hrsg. nach Vorlage}}}
\textit{dans toutes ses parties, et d'une force}
\edtext{\textit{egale à celle}}{%
\lemma{\textit{egale à}}%
\Bfootnote{%
\textit{(1)}~la puissance \textit{R}
\textit{(2)}~\textit{celle}~\textit{L}}}
\textit{de la puissance\indexS{puissance} \textit{R}}\lbrack,\rbrack\ \textit{on la peut regarder tant que dure cet equilibre,\indexS{equilibre}
comme divisée en autant de cordes; \textit{RRO}, \textit{RN}, et \textit{RM}
qu'il y a de poulies\indexS{poulie} dans l'echarpe \textit{LH}, et chacune de ces cordes\indexS{corde} comme fixe
en \textit{O}, \textit{M}, et \textit{N} et tirée du costé d'~\textit{R}, \textit{R}, \textit{R} 
par la puissance\indexS{puissance} \textit{R}.}
\textit{Et parce que les poulies\indexS{poulie} \textit{L}, \textit{K}, \textit{H} portent
chacune quelque chose du poids, regardons-le aussi comme divisé 
en autant de parties \textit{x}, \textit{y}, \textit{z} dont la partie \textit{x} soit portée
par la poulie\indexS{poulie} \textit{L}, et \textit{y} par \textit{K}, et \textit{z} par \textit{H}.}
\pend
\pstart
Or \textit{la proposition} des poulies\indexS{poulie}
$x:R \squaredots (a):b$\lbrack,\rbrack\
$y:R \squaredots (c):d$ et
$z:R \squaredots (e):f$\lbrack,\rbrack\
et poursuivant le calcul nous aurons
$D\ \text{(vel}\ x+y+z\text{)}:R \squaredots \overline{(a)df + (c)bf + (e)bd} : dfb$.\edlabel{37_05_217-218_3b}
\pend
\pstart
\edtext{%
Et seulement dans le cas où les cordes\indexS{corde} des poulies\indexS{poulie} \textit{qui touchent 
celles de d'escharpe \textit{LH} sont paralleles, la puissance\indexS{puissance} \textit{R} est au poids \textit{D},
comme $1$ au double du nombre des poulies\indexS{poulie} suspendues}\lbrack,\rbrack\
\edtext{\textit{mais} en}{\lemma{}\Bfootnote{\textit{mais} \bars{\textit{que} \textit{gestr.}} en~\textit{L}}}
\textit{tout autre cas elle luy est toujours en plus grande raison.}}{%
\lemma{et seulement \lbrack...\rbrack\ raison}%
\Cfootnote{\cite{02028}a.\,a.\,O., Coroll.~XVII, S.\;36 mit Änderungen.}}
\pend
\pstart
\edtext{%
Par \textit{cette proposition} generale on peut \textit{resoudre des problemes} comme \textit{par example}
\textit{faire soutenir quelque}
\edtext{\textit{poids} donné \textit{D}, \textit{à}}{%
\lemma{}%
\Bfootnote{\textit{poids} \bars{donné \textit{D} \textit{erg.}} 
\bars{\textit{D} \textit{streicht Hrsg.}} \textit{à}~\textit{L}}}
\textit{quelque puissance\indexS{puissance}} donnée \textit{\textit{R} par le moyen d'une seule poulie\indexS{poulie} à}
\edtext{\textit{moufle}. Et alors on peut}{%
\lemma{\textit{moufle}.}%
\Bfootnote{%
\textit{(1)}~On applique tousjours
\textit{(2)}~\bars{Et alors \textit{erg.}} on peut~\textit{L}}}
tousjours appliquer le moindre des deux au centre mobile de cette poulie\indexS{poulie}\indexS{centre mobile de la poulie} et le plus grand au bout libre 
\edtext{de la corde,\indexS{corde}}{%
\lemma{de la}%
\Bfootnote{%
\textit{(1)}~chorde
\textit{(2)}~corde,~\textit{L}}}%
}{%
\lemma{Par \lbrack...\rbrack\ corde}%
\Cfootnote{\cite{02028}a.\,a.\,O., S.\;36\textendash38 mit Änderungen.}}
\edtext{\textit{cependant} cela \textit{n'est pas}
\edtext{\textit{tousjours}}{\lemma{}\Bfootnote{\textit{tousjours} \textit{erg.~L}}}
\textit{necessaire, car si le plus grand de deux n'est pas plus que double de l'autre, 
on sera libre d'appliquer celuy qu'on voudra au centre mobile de la poulie\indexS{poulie}\indexS{centre mobile de la poulie} \textit{A}}\;\lbrack\!\!:\rbrack\
\textit{puisque qu'il y soit ou non}\lbrack,\rbrack\ \textit{il y peut tousjours monter 
jusqu'à estre double de l'autre,} et \textit{jamais davantage}.}{%
\lemma{\textit{cependant} \lbrack...\rbrack\ \textit{davantage}}%
\Cfootnote{\cite{02028}a.\,a.\,O., \textit{Remarque}, S.\;39 mit Änderungen.}}
\pend
\begin{edfigure}[0.5em]%
\includegraphics[scale=1]{%
gesamttex/edit_VIII-4/images/LH_37_05_217-218_d9_217v.pdf}
\hspace*{15mm}%
\includegraphics[scale=1]{%
gesamttex/edit_VIII-4/images/LH_37_05_217-218_d10_217v.pdf}%
\hspace*{5mm}%
\edcaption{%
\edtext{\lbrack\textit{Fig.~9}\rbrack}{%
\lemma{\lbrack\textit{Fig.~9}\rbrack}%
\Cfootnote[nonum]{\cite{02028}a.\,a.\,O., pl.~4, fig.~33 mit Änderungen.}}
\hspace*{48mm}%
\edtext{\edtext{\lbrack\textit{Fig.~10}\rbrack}{%
\lemma{\lbrack\textit{Fig.~10}\rbrack}%
\Cfootnote[nonum]{\cite{02028}a.\,a.\,O., pl.~4 fig.~33 mit Änderungen. Die Punkte $\omega$, $\beta$, $\kappa$, $\delta$ und die entsprechenden Verbindungsstrecken stammen von Leibniz. Ein gestrichener Entwurf zum Diagramm wird nicht wiedergegeben.}}}{%
\lemma{}%
\Afootnote[nonum]{\selectlanguage{latin}%
\textit{Neben und unterhalb dem Diagramm \lbrack\textit{Fig.~10}\rbrack}: %
Quid\textsuperscript{\lbrack a\rbrack}
si recta \textit{AD} non transeat per \textit{O}\lbrack?\rbrack\
Manifestum est idem oriri si pro \textit{A} sumas datum \textit{F} centrum gr.\indexS{centrum gravitatis}
et conatus\indexS{conatus} sumas esse $F\beta$, $F\kappa$, $F\delta$ quasi tunc contactus esset in $\omega$.
\newline
Potius conatus sustinentes\indexS{conatus sustinens} \textit{BA} et \textit{AD} in uno addendi 
qui component aequalem ipsi \textit{AC}, nam si plures sint superficies sustinentes\indexS{superficies sustinens} sic procedendum.
\\[1.5ex] %
{\footnotesize 
\textsuperscript{\lbrack a\rbrack}~Die verlorengegangenen Textteile dieser Anmerkung sind im S-Film Nr.~116 der GWLB Hannover noch vollständig lesbar.}}}%
\hspace*{15mm}}
\end{edfigure}
\pstart
\edtext{%
\edlabel{37_05_217-218_8a}%
\textit{Proposition fondamentale des poids soutenus sur quelque espece de surface\indexS{poids soutenus sur une surface} que ce soit
et pour toutes les directions possibles des puissances\indexS{puissance} qui y sont appliquées.
Telle que soit la surface \textit{GH}; le poids \textit{EO}, et la puissance\indexS{puissance} \textit{R} qui le soutient 
sont tousjours entre eux en raison reciproque des sinus des angles \indexS{sinus de l'angle}
que font leur lignes de direction\indexS{ligne de direction} avec \textit{AD} tirée perpendiculairement 
du point \textit{A} de leur concours sur la surface \textit{GH}.}\edlabel{37_05_217-218_8b}
Soit \textit{F} le centre de gravité\indexS{centre de gravité} du corps. Et la ligne de direction\indexS{ligne de direction} \textit{FC}.
\textit{1\textsuperscript{o}}\;Premièrement \textit{il faut qu'~\textit{FC} et \textit{EB} se rencontrent} 
comme \textit{en \textit{A}, car si elles estoient paralleles, le centre de gravité\indexS{centre de gravité} du poids 
pourroit encor descendre d'une longueur egale à} leur \textit{distance}.
\textit{2\textsuperscript{o}\;Les impressions sur \textit{A},}
\edtext{\lbrack\textit{estant}\rbrack}{\lemma{}\Bfootnote{sont \textit{L ändert Hrsg. nach Vorlage}}}
\textit{les memes que s'il estoit poussé par deus} puissances\indexS{puissance} \textit{suivant leur lignes de direction\indexS{ligne de direction}
\textit{AC} et \textit{AB}, le poids \textit{EO} ainsi soutenu}
\edtext{\lbrack\textit{sur}\rbrack}{\lemma{}\Bfootnote{par \textit{L ändert Hrsg. nach Vorlage}}}
\textit{la surface \textit{GH} par le concours}
\edtext{\textit{d'action de}}{\lemma{}\Bfootnote{\textit{d'action} \bars{\textit{de sa pesanteur} \textit{streicht Hrsg.}} \textit{de}~\textit{L}}}
\textit{sa pesanteur et de la puissance\indexS{puissance} \textit{R}, doit tendre suivant quelque ligne}
\edtext{\textit{AD}}{\lemma{}\Bfootnote{\textit{AD} \textit{erg.~L}}}
\textit{qui soit}
\edtext{\textit{diagonale}\indexS{diagonale du parallelogramme} du}{%
\lemma{\textit{diagonale}}%
\Bfootnote{%
\textit{(1)}~\textit{d'un par}
\textit{(2)}~du~\textit{L}}}
\textit{parallelogramme} \textit{BACD} et le poids et la puissance\indexS{puissance} sont entre eux 
\edtext{comme \textit{AB},}{%
\lemma{comme}%
\Bfootnote{%
\textit{(1)}~les
\textit{(2)}~\textit{AB},~\textit{L}}}
\textit{AC}.
\edtext{\textit{3\textsuperscript{o}}\;La \textit{ligne \textit{AD} doit estre perpendiculaire à la surface}}{%
\lemma{}%
\Afootnote{\selectlanguage{latin}%
\textit{Oberhalb der Zeile}: dubito an haec vera}}
\edtext{\lbrack\textit{HOG}\rbrack}{\lemma{}\Bfootnote{\textit{HOD} \textit{L ändert Hrsg.}}}
\edtext{\lbrack\textit{4\textsuperscript{o}}\rbrack\!}{\lemma{}\Bfootnote{4\textsuperscript{o} \textit{erg. Hrsg. nach Vorlage}}}
ou à sa touchante en \textit{O}, \textit{autrement} le corps \textit{roulerait} sur cette surface.
\edtext{\lbrack\textit{5\textsuperscript{o}}\rbrack\!}{\lemma{}\Bfootnote{4\textsuperscript{o} \textit{L ändert Hrsg. nach Vorlage}}}
\textit{R} est au poids \textit{F} \textit{comme le}
\edtext{\textit{sinus de l'angle\indexS{sinus de l'angle}}}{%
\lemma{\textit{sinus}}%
\Bfootnote{%
\textit{(1)}~\bars{de \textit{streicht Hrsg.}} l'angle \textit{BAD}, au sinus de l'angle \textit{EAD}
\textit{(2)}~\textit{de l'angle}~\textit{L}}}
\textit{\textit{DBA} ou de son complement\indexS{sinus du complement} \textit{BAC}, au sinus de l'angle\indexS{sinus de l'angle}
\textit{BDA} ou de \textit{DAC} qui luy est}
\edtext{\textit{egal.} Cette}{%
\lemma{\textit{egal}.}%
\Bfootnote{%
\textit{(1)}~\textit{Elle est aussi}
\textit{(2)}~Cette~\textit{L}}}
force \textit{est aussi pour la meme raison à la pesanteur de ce poids
comme le sinus du meme angle\indexS{sinus de l'angle} \textit{BAC}, au sinus de \textit{BAD}.
Ainsi la pesanteur de ce poids est à la force de cette puissance\indexS{puissance} comme le sinus
de l'angle\indexS{sinus de l'angle} \textit{BAD} au sinus de l'angle\indexS{sinus de l'angle} \textit{DAC},
et par consequent ce poids et cette puissance\indexS{puissance} sont entre eux 
en raison reciproque des sinus des angles\indexS{sinus de l'angle} que font leur lignes de direction\indexS{ligne de direction}
avec la ligne tirée de leur point de concours perpendiculairement
à la surface \textit{GH}, telle qu'elle soit.}}{%
\lemma{\textit{Proposition} \lbrack...\rbrack\ \textit{soit}}%
\Cfootnote{\cite{02028}a.\,a.\,O., S.\;40\,f. mit Änderungen.}}
\pend
\pstart
\href{\Scan{LH_37_05_218r}}%
{\lbrack218~r\textsuperscript{o}\rbrack}\ 
\edtext{%
On voit par là que si le corps
\edtext{\textit{EO}}{\lemma{}\Bfootnote{\textit{EO} \textit{erg.~L}}}
est une sphere\lbrack,\rbrack\ la ligne \textit{AD} passe necessairement 
par son centre de grandeur\indexS{centre de grandeur} (quand il ne seroit pas centre de gravité\indexS{centre de gravité})%
}{%
\lemma{On voit \lbrack...\rbrack\ gravité)}%
\Cfootnote{Siehe \cite{02028}a.\,a.\,O., Coroll.~III, S.\;43.}}
\selectlanguage{latin}%
(+~nam perpendicularis ad superficiem contactam \textit{HG}, 
est etiam perpendicularis ad superficiem sphaerae contingentem,
adeoque transit per centrum~\lbrack+\rbrack)
(+~si jam sphaera sit uniformis coincidunt \textit{A} et \textit{F}. 
Ergo linea \textit{BE} tendit etiam in centrum~+).
\edtext{%
Potentia\indexS{potentia} quantacunque sit pondus sustinere non potest in superficie,
nisi linea directionis\indexS{linea directionis} \textit{AR} intra angulum \textit{NRO}, 
qui est complementum anguli \textit{CAO} reperiatur.
Nam si \textit{AR} \textit{concouroit}
\edtext{(+~potius: estoit}{%
\lemma{\foreignlanguage{latin}{(+~potius:}}%
\Bfootnote{%
\textit{(1)}~esset coinc
\textit{(2)}~estoit~\textit{L}}}
coincidente~+)
\selectlanguage{french}%
avec 
\edtext{\textit{AN}, la puissance\indexS{puissance}}{%
\lemma{\textit{AN},}%
\Bfootnote{%
\textit{(1)}~\textit{R} poten
\textit{(2)}~la puissance~\textit{L}}}
\textit{R} \textit{seule porteroit tout le poids, et non avec} la surface \textit{contre} 
\edtext{\textit{l'hypothese. Si}}{%
\lemma{}%
\Bfootnote{\textit{l'hypothese.} \bars{Qui \textit{streicht Hrsg.}} \textit{Si}~\textit{L}}}
\textit{elle sortoit} hors \textit{de l'angle} (+~et venoit comme de (\textit{R})\textit{A}~+)
\textit{la diagonale\indexS{diagonale du parallelogramme} \textit{AD} du parallelogramme \textit{BC} se tourneroit vers \textit{G}}\lbrack,\rbrack\
\textit{ce qui feroit tomber le poids de ce costé là.}%
}{%
\lemma{Potentia \lbrack...\rbrack\ là}%
\Cfootnote{\cite{02028}a.\,a.\,O., Coroll.~IV, S.\;43 mit Änderungen.}}
\edtext{\textit{Le plan \textit{BAC} est aussi tousjours perpendiculaire 
à la surface que la ligne \textit{GH} represente,
et sur la quelle le poids est soutenu.}}{%
\lemma{\textit{Le plan} \lbrack...\rbrack\ \textit{soutenu}}%
\Cfootnote{\cite{02028}a.\,a.\,O., Coroll.~V, S.\;44 mit Auslassung.}}
\edtext{\textit{Le poids \textit{EO} est à la charge de} la \textit{surface} ou à sa resistence, 
\textit{comme le sinus de l'angle\indexS{sinus de l'angle} \textit{BAD} au sinus de l'angle\indexS{sinus de l'angle} \textit{BAC}.}}{%
\lemma{\textit{Le poids} \lbrack...\rbrack\ \textit{BAC}}%
\Cfootnote{\cite{02028}a.\,a.\,O., Coroll.~VI, S.\;44 mit Änderungen.}}
\edtext{Item \textit{le poids \textit{EO}}\lbrack,\rbrack\ {la charge (ou resistence) de la surface \textit{GH},
et la puissance\indexS{puissance} \textit{R}, sont entre eux comme}
\edtext{\textit{le}\lbrack\textit{s}\rbrack}{\lemma{}\Bfootnote{le \textit{L ändert Hrsg. nach Vorlage}}}
\textit{sinus des Angles\indexS{sinus de l'angle} \textit{BAD}, \textit{BAC} ou \textit{ABD} son complement,\indexS{sinus du complement}
et \textit{DAC} ou \textit{ADB} qui luy est egal, c'est à dire (lem.~5) comme les lignes
\textit{BD}, \textit{AD} et \textit{AB}.}}{%
\lemma{Item \lbrack...\rbrack\ \textit{AB}}%
\Cfootnote{\cite{02028}a.\,a.\,O., Coroll.~VII, S.\;44.}}
\edtext{\textit{Plus l'angle \textit{BAC} est obtus, moins la charge de la surface} sera \textit{grande}.}{%
\lemma{\textit{Plus} \lbrack...\rbrack\ \textit{grande}}%
\Cfootnote{\cite{02028}a.\,a.\,O., Coroll.~VIII, S.\;44 mit Änderung.}}
\edtext{\textit{Mais la charge de} la \textit{surface ne peut jamais estre plus grande que la somme} 
du \textit{poids et de} la \textit{puissance\indexS{puissance}} ensemble, lors qu'elle est 
\edtext{\lbrack\textit{horizontale}\rbrack}{\lemma{}\Bfootnote{perpendiculaire \textit{L ändert Hrsg. nach Vorlage}}}
\edtext{et les}{\lemma{}\Bfootnote{et \bars{toutes \textit{gestr.}} les~\textit{L}}}
deux puissances\indexS{puissance} tirent perpendiculairement.}{%
\lemma{\textit{Mais} \lbrack...\rbrack\ perpendiculairement}%
\Cfootnote{\cite{02028}a.\,a.\,O., Coroll.~IX, S.\;45 mit Änderungen.}}
\pend
\pstart
\edtext{\textit{On entend par} une \textit{\textls{surface horizontale}, un point de la surface courbe 
dont toutes les} touchantes (+~ou le plan touchant~+) \textit{sont horizontales}.}{%
\lemma{\textit{On entend} \lbrack...\rbrack\ \textit{horizontales}}%
\Cfootnote{\cite{02028}a.\,a.\,O., S.\;45 mit Änderungen.}}
\pend
\pstart
\edtext{Ce que nous disons d'un poids, et de sa direction qui est perpendiculaire à l'horizon,\indexS{horizon}
se peut entendre de toute autre puissance\indexS{puissance} quelque direction qu'elle ait.}{%
\lemma{Ce que \lbrack...\rbrack\ ait}%
\Cfootnote{Siehe \cite{02028}a.\,a.\,O., S.\;48.}}
\pend
\pstart
\edtext{%
Si des \textit{puissances\indexS{puissance}} diverses \textit{\textit{P} et \textit{R} soutiennent successivement
un meme poids \textit{A}, ou des poids egaux sur les points \textit{O} et \textit{Q}
d'un meme plan \textit{HG},} ces puissances\indexS{puissance} \textit{P} et \textit{R}, \textit{sont entre elles 
en raison composée de celles des sinus des angles\indexS{sinus de l'angle} \textit{QAD} et \textit{QAP}; 
\textit{OAR}, et \textit{OAD}.}}{%
\lemma{Si des \lbrack...\rbrack\ \textit{OAD}}%
\Cfootnote{\cite{02028}a.\,a.\,O., Coroll.~XXI, S.\;52 mit Änderungen.}}
\pend
\pstart
\edtext{On pourroit \textit{encor tirer} plusieurs \textit{corollaires par rapport aux poliplans\indexS{polyplan}}.}{%
\lemma{On pourroit \lbrack...\rbrack\ \textit{poliplans}}%
\Cfootnote{\cite{02028}a.\,a.\,O., S.\;55 mit Änderungen.}}
\selectlanguage{latin}%
(+~polyplans\indexS{polyplan} credo, id est cum multa plana simul sustinent.\indexS{multa plana simul sustinentia}
Et quando idem corpus a pluribus ejusdem superficiei punctis sustinetur,\indexS{corpus a pluribus ejusdem superficiei punctis sustentum}
idem esset ac si plures forent superficies. Et 
\edtext{semper utique}{%
\lemma{semper}%
\Bfootnote{%
\textit{(1)}~utrique
\textit{(2)}~utique~\textit{L}}}
sustentans considerari potest ut impellens. Illud vero hic considerandum quod
\edtext{semper pressio}{%
\lemma{semper}%
\Bfootnote{%
\textit{(1)}~debet press
\textit{(2)}~pressio~\textit{L}}}
superficiei et tensio chordae\indexS{chorda}\indexS{tensio chordae} simul, seu pressio plurium superficierum\indexS{pressio plurium superficierum} aequalis esse
toti ponderi videtur. Cujus vis\indexS{vis} nec perditur nec augetur.\indexS{vis nec perdita nec aucta} Quod explicabitur fingendo
pro iis pondera elevata. Sed hoc intelligendum arbitror non de pondere\indexS{pondus} sed de momento.\indexS{momentum}
Res accurate consideranda. 
\edtext{Nam demonstravi 
tunc pondus non amplius descendere, 
quia pondera chordis appensa\indexS{pondus chordae appensum} magis ascendant, cum chordae\indexS{chorda}
ita dispositae,
ut ostendi in scheda separata.}{%
\lemma{Nam \lbrack...\rbrack\ separata}%
\Cfootnote{Die Anspielung konnte noch nicht ermittelt werden.}}
\pend
\begin{edfigure}[1em][1em]
\includegraphics[scale=1]{%
gesamttex/edit_VIII-4/images/LH_37_05_217-218_d11_218r.pdf}
\edcaption{\edtext{\lbrack\textit{Fig.~11}\rbrack}{%
\lemma{\lbrack\textit{Fig.~11}\rbrack}%
\Cfootnote[nonum]{Ein gestrichener Entwurf zum Diagramm wird nicht wiedergegeben.}}}
\end{edfigure}
\selectlanguage{latin}
\pstart
Quid si lineae perpendiculares ad superficiem tangentem non concurrant in linea directionis\indexS{linea directionis}\lbrack?\rbrack\
Considerandum 
\edtext{est pondus}{%
\lemma{est}%
\Bfootnote{%
\textit{(1)}~lineam
\textit{(2)}~pondus~\textit{L}}}
\textit{F} velut actum circa centrum
\edtext{\textit{E}}{\lemma{}\Bfootnote{\textit{E} \textit{erg.~L}}}
agere in \textit{O}, et vicissim pondus \textit{F} velut actum circa centrum \textit{O}, agere in \textit{E}.
Cogitandum 
\edtext{est plana sustentantia\indexS{planum sustentans} esse ut lineas}{%
\lemma{est}%
\Bfootnote{%
\textit{(1)}~duas lineas esse sustentatas
\textit{(2)}~plana sustentantia
\textit{(a)}~cog
\textit{(b)}~esse ut
\textit{(aa)}~rectas
\textit{(bb)}~lineas~\textit{L}}}
rigidas \textit{TOI}, \textit{LEK} sustentatas in duobus punctis \textit{T} et \textit{I}, vel 
\edtext{\lbrack\textit{L} et \textit{K}\rbrack.}{\lemma{}\Bfootnote{\textit{LK} \textit{L ändert Hrsg.}}}
Vel quod eodem redit pali ad superficiem
\edtext{corpori}{\lemma{}\Bfootnote{corpori \textit{erg.~L}}}
perpendiculares \textit{ED}, \textit{OQ}. Quid 
\edtext{si superficies}{%
\lemma{si}%
\Bfootnote{%
\textit{(1)}~in planis
\textit{(2)}~superficies~\textit{L}}}
sustinens\indexS{superficies sustinens} et 
\edtext{sustentata\indexS{superficies sustentata} in}{%
\lemma{sustentata}%
\Bfootnote{%
\textit{(1)}~seu
\textit{(2)}~in~\textit{L}}}
linea vel plano tangant, quae puncta eligemus?
Hinc videtur principium difficultat\textlangle is\textrangle\
\textlangle sol\textrangle vendae praeberi,
semper enim fingi potest salva sustentandi ratione,
in plano ejus magno esse contactus, sic ut pa\textlangle li\textrangle\ 
e\textlangle li\textrangle g\textlangle an\textrangle tur,
qui producti incidant in centrum gravitatis\indexS{centrum gravitatis} corporis.
NB NB. Per figuram appictam res accurate poterit determinari. 
Ope ponderum sumamus $VQ = V\hebr{.h}$ et $ND = N\hebr{l}$. 
Sed haec alias.~+)
\pend
\begin{edfigure}[0.25em][1em]
\includegraphics[scale=1]{%
gesamttex/edit_VIII-4/images/LH_37_05_217-218_d12_218r.pdf}
\edcaption{\edtext{\lbrack\textit{Fig.~12}\rbrack}{%
\lemma{\lbrack\textit{Fig.~12}\rbrack}%
\Cfootnote[nonum]{\cite{02028}a.\,a.\,O., pl.~6, fig.~42.}}}
\end{edfigure}
\selectlanguage{french}
\pstart
\edtext{%
\textit{\textls{Proposition fondamentale} pour toutes sortes des leviers,\indexS{levier}
de quelque espece et dans quelque situation qu'ils soyent et pour toutes
les directions possibles des puissances\indexS{puissance} ou des poids qui y sont appliqués.}
\textit{Soyent les puissances\indexS{puissance} \textit{E} et \textit{F} appliqués aux points \textit{O} et \textit{X}
du levier\indexS{levier} \textit{MN} de quelque espece et en quelque situation qu'il soit, 
quelque angle \textit{OAX} que fassent aussi entre elles les lignes de direction\indexS{ligne de direction}
de ces puissances\indexS{puissance}}
\edtext{\textit{indefi}\lbrack\textit{ni}\rbrack\textit{ment}}{\lemma{}\Bfootnote{indefiment \textit{L ändert Hrsg. nach Vorlage}}}
\textit{prolongées vers \textit{A}}\;\lbrack\!\!:\rbrack\ 
\textit{ces deux puissances\indexS{puissance} feront equilibre sur le point fixe \textit{B} de ce levier,\indexS{levier}
par ou passe la diagonale\indexS{diagonale du parallelogramme} \textit{AG} du parallelogramme \textit{RS}, dont les costés
\textit{AS} et \textit{AR} sont entre eux comme les puissances\indexS{puissance} \textit{E} et \textit{F}.}
\textit{Concevons pour un moment que \textit{MAN} est la figure de ce levier\indexS{levier}}
\foreignlanguage{latin}{(+~probandum an hoc licere~+)}
\textit{et qu'au lieu d'estre tiré ou poussé par les puissances\indexS{puissance} \textit{E} et \textit{F}
son point \textit{A} soit seulement poussé, mais en meme temps vers \textit{R} suivant \textit{AR}}
\edtext{\textit{par une puissance\indexS{puissance} egale à \textit{F}}}{\lemma{}\Bfootnote{\textit{par une puissance egale à \textit{F}} \textit{erg.~L}}}
\textit{et vers \textit{S} suivant \textit{AS} par une puissance\indexS{puissance} egale à \textit{E}. Il est constant que
l'impression composée\indexS{impression composée} que ce point recevroit alors du concours d'action 
de ces deux puissances\indexS{puissance}}
ainsi ce point \textit{tendra}
\edtext{\textit{en \textit{AG}}, et}{%
\lemma{\textit{en \textit{AG}},}%
\Bfootnote{%
\textit{(1)}~et soutiendra
\textit{(2)}~et~\textit{L}}}
\textit{le point fixe} comme \textit{B} en \textit{GAB} 
\textit{soutiendra ces puissances,\indexS{puissance} et par consequent le levier\indexS{levier}}
\edtext{\lbrack\textit{n'}\rbrack\textit{ayant}}{\lemma{}\Bfootnote{\textit{ayant} \textit{L ändert Hrsg. nach Vorlage}}}
\textit{que ce que} ces puissances\indexS{puissance} luy donnent sera \textit{en equilibre.}\indexS{equilibre}
\textit{Remettons les deux premieres \textit{E} et \textit{F},} elles seront le 
\edtext{meme sur}{%
\lemma{meme}%
\Bfootnote{%
\textit{(1)}~et on peut
\textit{(2)}~sur~\textit{L}}}
\textit{MAN}, \textit{et si on en retranche \textit{MAN},} et laisse seulement \textit{MN},
ce sera le meme.}{%
\lemma{\textit{\textls{Proposition}} \lbrack...\rbrack\ meme}%
\Cfootnote{\cite{02028}a.\,a.\,O., S.\;58\,f. mit Änderungen.}}
\selectlanguage{latin}%
(+~Non admodum placet argumentandi
\edlabel{37_05_217-218_6a}%
\edtext{}{%
\xxref{37_05_217-218_6a}{37_05_217-218_6b}%
\lemma{ratio.~+)}%
\Bfootnote{%
\textit{(1)}~Donc le point
\textit{(2)}~Ex his apparet,~\textit{L}}}%
ratio.~+) 
\edtext{Ex his apparet,\edlabel{37_05_217-218_6b}
\selectlanguage{french}%
encor comme le point fixe du levier\indexS{point fixe du levier}\indexS{levier} est chargé sçavoir par la puissance\indexS{puissance} comme \textit{AG}.}{%
\lemma{Ex his \lbrack...\rbrack\ \textit{AG}}%
\Cfootnote{Siehe \cite{02028}a.\,a.\,O., Coroll.~IV, S.\;61.}}
\edtext{\textit{Personne que je sçache n'avoit encor demonstré la charge ou les points d'appui des leviers.\indexS{levier}\indexS{point d'appui du levier}
Il y a bien des problemes qu'on ne sçavoit resoudre sans cela, comme par exemple:
quelles doivent estre les directions de deux puissances\indexS{puissance} soit egales soit inegales
pour qu'elles fassent equilibre\indexS{equilibre} sur quelque levier\indexS{levier} que ce soit, dont l'appui est une sphère.}
Item \textit{sur combien de points} le \textit{levier\indexS{levier} ainsi appuyé, il est possible qu'elles fassent
equilibre\indexS{equilibre} en changeant seulement leur} 
\edtext{\textit{directions.} Item}{%
\lemma{\textit{directions.}}%
\Bfootnote{%
\textit{(1)}~\bars{\textit{Il} \textit{streicht Hrsg.}} \textit{n'e}
\textit{(2)}~Item~\textit{L}}}
\textit{trouver le point d'appui\indexS{point d'appui du levier} de} 
\edtext{\textit{celuy}}{\lemma{}\Afootnote{\foreignlanguage{latin}{\textit{Oberhalb celuy}:\enskip levier}}}
\textit{au quel tant de puissances\indexS{puissance} qu'on voudra soyent appliquées pour toutes les directions possibles.}
Item \textit{deux puissances\indexS{puissance} estant} 
\edtext{\textit{données avec}}{%
\lemma{\textit{données}}%
\Bfootnote{%
\textit{(1)}~\textit{trouver} nat
\textit{(2)}~hinc
\textit{(3)}~\textit{avec}~\textit{L}}}
\textit{leur directions et leur points} d'appui
\edtext{de}{\lemma{}\Bfootnote{de \textit{erg.~L}}}
\textit{trouver quelle doit estre la direction et le point d'application d'une 3\textsuperscript{me} puissance\indexS{puissance}
aussi donnée pour que toutes les trois ensemble fassent equilibre\indexS{equilibre} 
sur quelque point donné que ce soit de ce levier.\indexS{levier}}%
}{%
\lemma{\textit{Personne} \lbrack...\rbrack\ \textit{levier}}%
\Cfootnote{\cite{02028}a.\,a.\,O., S.\;64 mit Änderungen.}}
\edtext{\textit{On en voit} aussi, \textit{que lors que deux puissances\indexS{puissance} font equilibre sur un levier,\indexS{levier} elles sont tousjours
entre elles en raison reciproque des distances de leur lignes de direction,\indexS{ligne de direction} au point d'appui de ce levier.\indexS{levier}\indexS{point d'appui du levier}
Et} viceversa quand \textit{elles sont en raison reciproque de} ces \textit{distances, elles font equilibre.\indexS{equilibre}
Voilà ce qu'on appelle ordinairement le premier principe de Mecanique.\indexS{premier principe de Mecanique}}}{%
\lemma{\textit{On} \lbrack...\rbrack\ \textit{Mecanique}}%
\Cfootnote{\cite{02028}a.\,a.\,O., Coroll.~XIII, S.\;66 mit Änderungen.}}
\selectlanguage{latin}%
(+~a solidis seu rigidis merito incipiunt potius quam a 
\edtext{chordis.\indexS{chorda}~\lbrack+)\rbrack\ 
\selectlanguage{french}%
\textit{Lemme}}{%
\lemma{chordis.}%
\Bfootnote{%
\textit{(1)}~\textit{Si deux}
\textit{(2)}~\textit{Lemme}~\textit{L}}}
\edtext{\textit{6. Si deux puissances\indexS{puissance} \textit{D} et \textit{E} appliqués au meme point \textit{A},
d'un levier\indexS{levier} \textit{AB} dont le point fixe\indexS{point fixe du levier} soit \textit{B}}\lbrack,\rbrack\ \textit{sont entre elles en raison reciproque 
des sinus des angles\indexS{sinus de l'angle} que font leur lignes de direction\indexS{ligne de direction} avec ce levier\indexS{levier}}\lbrack,\rbrack\
\textit{elles demeureront en equilibre.\indexS{equilibre}}}{%
\lemma{\textit{6.~Si} \lbrack...\rbrack\ \textit{equilibre}}%
\Cfootnote{\cite{02028}a.\,a.\,O., Lemme~VI, S.\;66\,f.}}
La demonstration est telle:
\pend
\pstart\noindent
\href{\Scan{LH_37_05_218v}}%
{\lbrack218~v\textsuperscript{o}\rbrack}\ 
\pend
\nobreak
\begin{edfigure}[0em]%
\hspace{7mm}%
\includegraphics[scale=1]{%
gesamttex/edit_VIII-4/images/LH_37_05_217-218_d13_218v.pdf}
\edcaption{\edtext{\lbrack\textit{Fig.~13}\rbrack}{%
\lemma{\lbrack\textit{Fig.~13}\rbrack}%
\Cfootnote[nonum]{\cite{02028}a.\,a.\,O., pl.~7, fig.~47.}}}
\end{edfigure}
\pstart
\edtext{%
\textit{Le point \textit{A} ainsi tiré par les puissances\indexS{puissance} \textit{D} et \textit{E}, doit en recevoir 
une impression suivant quelque ligne \textit{AG} diagonale\indexS{diagonale du parallelogramme} d'un parallelogramme \textit{RS},
fait sous des parties de leur lignes de direction,\indexS{ligne de direction} \textit{AS}, \textit{AR}, 
qui soyent entre elles (Lemme~3.) comme ces puissances\indexS{puissance}}\,\lbrack\!;\rbrack\
\textit{et la force de cette impression doit estre} 
\edtext{\lbrack(\rbrack\textit{Lem.~3}\lbrack)\rbrack}{\lemma{}\Bfootnote{\textit{Lem.~3} \textit{erg.~L}}}
\textit{a chacune des puissances\indexS{puissance} \textit{D} et \textit{E}}
\edtext{\textit{comme \textit{AG}}}{%
\lemma{}%
\Bfootnote{%
\textit{comme} \bars{comme \textit{streicht Hrsg.}} \textit{AG}~\textit{L}}}
\textit{à chacune des lignes \textit{AS}, et \textit{AR} ou \textit{SQ} qui luy est egale,
c'est à dire (Lem.~5) comme le sinus de l'angle\indexS{} \textit{ASG} à chacun des sinus 
des angles\indexS{} \textit{AGS} ou \textit{GAR} qui lui est egal, et \textit{GAS}, 
et par consequent les puissances\indexS{puissance} \textit{D} et \textit{E} sont entre elles
comme les sinus \textit{GAR}, et \textit{GAS}. Elles sont aussi (hyp.) entre elles
comme les sinus des Angles\indexS{} \textit{RAB} et \textit{SAB}, et par consequent
les sinus de \textit{GAR}, et \textit{RAB}, de meme que ceux de \textit{GAS}, et \textit{SAB},
sont egaux entre eux. Donc} leur \textit{angles sont aussi egaux, ou dumoins complemens
l'un de l'autre à deux droits. D'où il suit qu'\textit{AG} est en ligne droite avec \textit{AB},
et par consequent le point fixe\indexS{point fixe du levier} \textit{B} du levier\indexS{levier} \textit{AB}, soutient luy seul
l'impression toute entiere que le point \textit{A} reçoit du concours d'action}
\textit{des puissances\indexS{puissance} \textit{D} et \textit{E}, et ce point demeurant ainsi en repos}\lbrack,\rbrack\
\textit{ces puissances\indexS{puissance} doivent aussi demeurer par consequent en equilibre\indexS{equilibre}}\lbrack,\rbrack\
\textit{ce qu'il falloit demonstrer.}}{%
\lemma{\textit{Le point} \lbrack...\rbrack\ \textit{demonstrer}}%
\Cfootnote{\cite{02028}a.\,a.\,O., Demonstration, S.\;67 mit Änderungen.}}
\pend
\pstart
\foreignlanguage{latin}{Sequitur}
\edtext{%
\textit{Autre proposition des leviers.\indexS{levier}
Si les puissances\indexS{puissance} \textit{E} et \textit{F} appliquées aux differens points \textit{X} et \textit{O}
du levier\indexS{levier} \textit{BO} sont entre elles en raison reciproque des distances \textit{BD} et \textit{BP}
de leurs lignes de direction\indexS{ligne de direction} \textit{XE}, et \textit{OF}, qu'elles quelles soyent, au point fixe\indexS{point fixe du levier}
de ce levier,\indexS{levier} c'est à dire si $E : F\squaredots BP : BD$}\lbrack,\rbrack\
\textit{ces deux puissances\indexS{puissance} ainsi appliquées demeurent en equilibre\indexS{equilibre} sur ce levier.\indexS{levier}}}{%
\lemma{\textit{Autre} \lbrack...\rbrack\ \textit{levier}}%
\Cfootnote{\cite{02028}a.\,a.\,O., S.\;70 mit Änderungen.}}
\edtext{%
\textit{Demonstration. Si la puissance\indexS{puissance} \textit{E} étoit appliquée en \textit{O} 
suivant une ligne de direction\indexS{ligne de direction} \textit{OQ}, dont la distance \textit{BQ} au point d'appui\indexS{point d'appui du levier} \textit{B}
de ce levier,\indexS{levier} fut egale a \textit{BD}, qui est celle de ce meme point à la ligne de direction\indexS{ligne de direction} \textit{XE}
qu'elle a presentement, cette puissance\indexS{puissance} (Lemme~6.) demeureroit en equilibre\indexS{equilibre} 
avec la puissance\indexS{puissance} \textit{F}. Or l'action de cette puissance\indexS{puissance} appliquée en \textit{X}, 
suivant \textit{XE} est la meme}
\edtext{\textit{(Lem.~6) qu'elle}}{%
\lemma{}%
\Bfootnote{%
\textit{(Lem.~6} \bars{\textit{cor.}~\textit{gestr.}} \textit{qu'elle}~\textit{L}}}
\textit{seroit alors sur ce levier.\indexS{levier} Donc ces puissances\indexS{puissance} ainsi appliquées en \textit{X} et en \textit{O}
doivent demeurer en equilibre\indexS{equilibre}}\lbrack,\rbrack\ \textit{ce qu'il falloit demonstrer.}}{%
\lemma{\textit{Demonstration} \lbrack...\rbrack\ \textit{demonstrer}}%
\Cfootnote{\cite{02028}a.\,a.\,O., Demonstration, S.\;70\,f. mit Änderung.}}
\pend
\pstart
\edtext{Ces raisonnemens se peuvent appliquer 
\textit{à la balance}\lbrack,\rbrack\ \textit{à la romaine, au tour,
aux roues à dent, aux ciseaux, aux tenailles, aux etocs, etc.}}{%
\lemma{Ces \lbrack...\rbrack\ \textit{etc.}}%
\Cfootnote{\cite{02028}a.\,a.\,O., S.\;71 mit Änderungen.}}
\pend
\begin{edfigure}[1em][1em]
\includegraphics[scale=1]{%
gesamttex/edit_VIII-4/images/LH_37_05_217-218_d14_218v.pdf}
\edcaption{%
\edtext{\lbrack\textit{Fig.~14}\rbrack}{%
\lemma{\lbrack\textit{Fig.~14}\rbrack}%
\Cfootnote[nonum]{\cite{02028}a.\,a.\,O., pl.~8, fig.~56. Teilverlust wegen Papierschadens; die Buchstaben \textit{E}, \textit{G}, \textit{H}, \textit{b}, \textit{e} und \textit{h} ergänzt Hrsg. nach der Vorlage.
Ein gestrichener Entwurf zum Diagramm wird nicht wiedergegeben.}}}
\end{edfigure}
\pstart
\edtext{%
\textit{Ce que nous venons de demonstrer des leviers droits\indexS{levier droit} se peut appliquer aux courbes,\indexS{levier courbe}
en faisant à discretion quelque ligne droite \textit{AH}, qui rencontre il n'importe comment
toutes les lignes de direction\indexS{ligne de direction} des puissances\indexS{puissance} \textit{M}, \textit{N}, \textit{O}, \textit{P}, \textit{Q}
en \textit{a}, \textit{c}, \textit{e}, \textit{h}, \textit{g}}\lbrack,\rbrack\ 
\textit{car la regardant comme un levier\indexS{levier}} ou \textit{ces puissances\indexS{puissance} soyent appliquées, 
on trouvera le point d'appui\indexS{point d'appui du levier} \textit{b}, avec sa ligne de direction\indexS{ligne de direction} \textit{bB}, 
qui rencontrant le levier\indexS{levier} \textit{AH} donnera le point \textit{B}, qui sera son point d'appui.\indexS{point d'appui du levier}}}{%
\lemma{\textit{Ce} \lbrack...\rbrack\ \textit{appui}}%
\Cfootnote{\cite{02028}a.\,a.\,O., Coroll.~II, S.\;73 mit Änderungen.}}
\selectlanguage{latin}%
Quod si directiones non parallelae\lbrack?\rbrack\
quaeratur punctum concursus, 
\edtext{et facto}{%
\lemma{et}%
\Bfootnote{%
\textit{(1)}~ducto
\textit{(2)}~facto~\textit{L}}}
parallelogrammo, habebitur duarum potentiarum\indexS{potentia} concursus,\indexS{concursus duarum potentiarum} unde et trium\indexS{concursus trium potentiarum} etc.
\pend
\selectlanguage{french}
\pstart
\edtext{La vis n'est qu'une espece de plan incliné\indexS{plan incliné} ou surface. En voici une proposition.}{%
\lemma{La vis \lbrack...\rbrack\ proposition}%
\Cfootnote{\cite{02028}a.\,a.\,O., S.\;77.}}
\edtext{%
\textit{Lors qu'une puissance\indexS{puissance} soutient quelque poids, ou l'action de quelque autre force
à l'aide d'une vis}\lbrack,\rbrack\ \textit{soit que cette vis soit fixe ou que ce soit son ecroue}\lbrack,\rbrack\
\textit{cette puissance\indexS{puissance} est tousjours à ce poids ou à cette force\lbrack,\rbrack\ quelle qu'elle soit\lbrack,\rbrack\
comme la distance qui est}
\edtext{\textit{entre deux}}{%
\lemma{\textit{entre}}%
\Bfootnote{%
\textit{(1)}~cette vis
\textit{(2)}~\textit{deux}~\textit{L}}}
\textit{des pas de cette vis à la circomference d'un cercle dont le rayon est egal à la distance
qui est entre cette puissance\indexS{puissance} et l'axe de cette meme vis.}}{%
\lemma{\textit{Lors} \lbrack...\rbrack\ \textit{vis}}%
\Cfootnote{\cite{02028}a.\,a.\,O., S.\;79\,f.}}%
\edlabel{37_05_217-218_4b}		%
\pend
\begin{edfigure}%
\includegraphics[scale=1]{%
gesamttex/edit_VIII-4/images/LH_37_05_217-218_d15_218v.pdf}%
\hspace*{7mm}
\edcaption[0.25em]{\edtext{\lbrack\textit{Fig.~15}\rbrack}{%
\lemma{\lbrack\textit{Fig.~15}\rbrack}%
\Cfootnote[nonum]{%
\cite{02028}a.\,a.\,O., pl.~1, fig.~1;
die vier zum \textit{Examen} gehörigen Tafeln befinden sich am Ende des Bandes.
Zwei gestrichene Entwürfe zum Diagramm werden nicht wiedergegeben.}%
\lemma{}%
\Afootnote[nonum]{%
\selectlanguage{ngerman}%
\textit{Neben dem Diagramm \lbrack\textit{Fig.~15}\rbrack}:
\textit{CP} et \textit{CQ} 
\indexN{\textls{Varignon}, Pierre de 1654\textendash1722}l'auteur 
appelle sublimités des puissances.\indexS{puissance}  
Si elles estoient au dessous de \textit{C},\textsuperscript{\lbrack a\rbrack} 
il les appelleroit \textls{profondeurs},\textsuperscript{\lbrack b\rbrack} 
et \textit{CN}, \textit{CK} \textls{parties proportionelles}.\textsuperscript{\lbrack c\rbrack}
\\[1.5ex] %
{\footnotesize
\textsuperscript{\lbrack a\rbrack}~\textit{C}, \textit{(1)}~seroient \textit{(2)}~il~\textit{L}\quad 
\textsuperscript{\lbrack b\rbrack}~\textls{profondeurs}, \textit{(1)}~la somme \textit{(2)}~et~\textit{L}\quad
\textsuperscript{\lbrack c\rbrack}~\textit{CP} \lbrack...\rbrack\ \textls{proportionelles}:
\indexN{\textls{Varignon}, Pierre de 1654\textendash1722}\textsc{Varignon},
\cite{02028}\textit{Examen}, S.\;95 und 100.}}}}
\end{edfigure}
\pstart
\edlabel{37_05_217-218_5a}%
\edtext{%
\indexN{\textls{Borelli} (Borellus), Giovanni Alfonso 1608\textendash1679}M.\;Borelli
tome I \cite{02099}\textit{De mot. animal.} prop.~68 %
monstre que 
\textit{lors que les puissances\indexS{puissance} \textit{R.S} soutiennent ensemble le poids \textit{T}, la}
\edtext{\textit{puissance\indexS{puissance} \textit{R}}}{%
\lemma{}%
\Bfootnote{%
\textit{puissance} \bars{puissance \textit{streicht Hrsg.}} \textit{R}~\textit{L}}}
\textit{soutient pour sa part une partie \textit{X} de meme qu'elle feroit si elle estoit appliquée}
\edtext{\textit{suivant sa}}{%
\lemma{\textit{suivant}}%
\Bfootnote{%
\textit{(1)}~cette
\textit{(2)}~\textit{sa}~\textit{L}}}
\textit{meme direction \textit{AC} avec cette partie \textit{X} au levier\indexS{levier} horizontal \textit{CG},
et la puissance\indexS{puissance} \textit{S}, soutient \textit{Z},} comme \textit{si elle estoit appliquée
au levier\indexS{levier} horizontal \textit{CI} egal au premier. 
Et par consequent si} selon notre \textit{lemme 3 on regarde l'impression\indexS{impression composée} de \textit{S} sur \textit{C},
comme composée des impressions \textit{CO} et \textit{CH}, on trouvera que}
l'impression selon \textit{CO}, \textit{est egal à la resistence du levier\indexS{levier} \textit{CG}
pour empecher la corde\indexS{corde} \textit{ACX} de se redresser, c'est à dire 
egal à la charge de l'appui \textit{G} de ce levier.\indexS{levier} Or} 
par la \textit{prop. fondamentale des leviers\indexS{levier} la puissance\indexS{puissance} \textit{R} est à la charge de cet appui comme}
\edtext{\textit{\textit{AC}, costé}}{%
\lemma{\textit{AC},}%
\Bfootnote{%
\textit{(1)}~à \textit{CG}
\textit{(2)}~\textit{costé}~\textit{L}}}
\textit{du parallelogramme \textit{AE}, à la diagonale\indexS{diagonale du parallelogramme} \textit{CG}, et} de meme
\textit{l'impression de \textit{S} sur \textit{C}, suivant \textit{CO},}
\edtext{\textit{est à} l'impression}{%
\lemma{\textit{est à}}%
\Bfootnote{%
\textit{(1)}~\textit{B}, comme
\textit{(2)}~l'impression~\textit{L}}}
de \textit{S} \textit{suivant \textit{CB} c'est à dire à cette puissance\indexS{puissance} elle meme,}
\edtext{\lbrack\textit{comme}\rbrack}{\lemma{}\Bfootnote{\textit{comme} \textit{erg. Hrsg. nach Vorlage}}}
\textit{\textit{OC} costé du parallelogramme \textit{OH} à sa diagonale\indexS{diagonale du parallelogramme} \textit{BC}, c'est}
\edtext{\textit{à dire (en}}{%
\lemma{\textit{à dire}}%
\Bfootnote{%
\textit{(1)}~comme
\textit{(2)}~\textit{(en}~\textit{L}}}
\textit{faisant \textit{IK} perpendiculaire sur \textit{BC}}\lbrack)\rbrack,
\textit{comme \textit{CK}, à \textit{CI} egale (\textup{per} hyp.) à \textit{CG}, donc}
\edtext{\lbrack\textit{la}\rbrack}{\lemma{}\Bfootnote{\textit{la} \textit{erg. Hrsg. nach Vorlage}}}
\textit{puissance\indexS{puissance} \textit{R} est à \textit{S}, comme le produit
de \textit{AC} par \textit{CK} est au quarré de \textit{CG}. Or en faisant \textit{GF}
perpendiculaire sur \textit{AC} les triangles \textit{AGC} et \textit{GFC} estant semblables}\lbrack,\rbrack\
\textit{le produit de \textit{AC} par \textit{CF}, est egal au quarré de \textit{CG}}\;\lbrack\!\!:\rbrack\
\textit{donc la puissance\indexS{puissance} \textit{R} est à la puissance\indexS{puissance} \textit{S} comme le produit de \textit{AC} par \textit{CK} 
au produit de la meme \textit{AC} par \textit{CF}, c'est à dire comme \textit{CK} à \textit{CF}, 
ou bien à cause des}
\edtext{\textit{ra}\lbrack\textit{y}\rbrack\textit{ons}}{%
\lemma{}\Bfootnote{raisons \textit{L ändert Hrsg. nach Vorlage}}}
\textit{\textit{IC}, \textit{CG} (hyp.) egaux\lbrack,\rbrack\ comme le sinus de \textit{KIC} egal à \textit{BCH} 
au sinus de \textit{FGC}, egal à \textit{ACH}. Donc la puissance\indexS{puissance} \textit{R} est à la puissance\indexS{puissance} \textit{S} 
comme le sinus de l'angle\indexS{sinus de l'angle} \textit{BCH} à celuy de \textit{ACH}, 
qui a cause du parallelogramme \textit{MN} est egal à \textit{CHM}, donc (lemme~5) la puissance\indexS{puissance}} 
\edtext{\lbrack\textit{R}\rbrack}{%
\lemma{}%
\Bfootnote{\textit{S} \textit{L ändert Hrsg. nach Vorlage}}} 
\textit{est à la puissance\indexS{puissance}} 
\edtext{\lbrack\textit{S}\rbrack}{%
\lemma{}%
\Bfootnote{\textit{R} \textit{L ändert Hrsg. nach Vorlage}}} 
\textit{comme \textit{HM}, ou comme son egale \textit{NC}, à \textit{CM}.}}{%
\lemma{M.\;Borelli \lbrack...\rbrack\ à \textit{CM}}%
\Cfootnote{\cite{02028}a.\,a.\,O., S.\;93\textendash95 mit Änderungen. 
Vgl. \indexN{\textls{Borelli} (Borellus), Giovanni Alfonso 1608\textendash1679}\textsc{G.\,A.~Borelli},
\cite{02099}\textit{De motu animalium}, Bd~I, Rom 1680, prop.~LXVIII mit Scholium, S.\;129\,f.}}
\edtext{%
Cela prouvé\lbrack,\rbrack\ 
\indexN{\textls{Borelli} (Borellus), Giovanni Alfonso 1608\textendash1679}M.\;Borelli 
pouvoit passer plus avant et 
\textit{faisant \textit{MP} et \textit{NQ} perpendiculaires sur \textit{HC}, ces lignes
marquant tousjours \textit{CP} egale à \textit{HQ}}\lbrack,\rbrack\ 
\textit{cet auteur auroit trouvé comme il a fait p.~137.} edit. de Rome
\textit{que chacune des puissances\indexS{puissance} \textit{R} et \textit{S} estant tousjours 
par le coroll. de sa 69. Prop. à tout ce poids comme chacun des costés
\textit{CN}, \textit{CM} du parallelogr. \textit{MN}, à la somme de leur sublimités
\textit{CP}, et \textit{CQ}, lui est aussi tousjours comme chacun de ces memes costés
à la diagonale\indexS{diagonale du parallelogramme} \textit{CH} de ce meme parallelogramme.} 
\selectlanguage{latin}%
Q.\,E.\,Dem.}{%
\lemma{Cela \lbrack...\rbrack\ Q.\,E.\,Dem.}%
\Cfootnote{\indexN{\textls{Varignon}, Pierre de 1654\textendash1722}\textsc{Varignon},
\cite{02028}\textit{Examen}, S.\;95 mit Änderungen.
Vgl. \indexN{\textls{Borelli} (Borellus), Giovanni Alfonso 1608\textendash1679}\textsc{Borelli},
\cite{02099}\textit{De motu animalium}, prop.~LXIX mit Corollarium und Digressio, S.\;131\textendash143, hier S.\;137.}}
\edtext{Hinc concludere potuisset, non posse mutare angulos quin mutatur aequilibrium.\indexS{aequilibrium}}{%
\lemma{Hinc \lbrack...\rbrack\ aequilibrium}%
\Cfootnote{\indexN{\textls{Varignon}, Pierre de 1654\textendash1722}\textsc{Varignon},
\cite{02028}\textit{Examen}, S.\;96.}}
Contra 
\edtext{demonstrationem, (quae \indexN{\textls{Pardies}, Ignace Gaston 1636\textendash1673}Pardiesii videtur)}{%
\lemma{demonstrationem, (quae Pardiesii videtur)}%
\Cfootnote{\cite{02028}a.\,a.\,O., S.\;97.
\indexN{\textls{Borelli} (Borellus), Giovanni Alfonso 1608\textendash1679}Borelli schreibt den Beweis 
lediglich einem \glqq insignis geometra neotericus\grqq\ zu (\cite{02099}\textit{De motu animalium}, S.\;135\,f.); Varignon kommentiert, dass der Beweis \glqq paroît être du P.~Pardie\grqq.
Vgl. \indexN{\textls{Pardies}, Ignace Gaston 1636\textendash1673}\textsc{I.\,G.~Pardies}, \textit{La statique}, Paris 1673, prop.~LXVII, S.\;112\textendash117.}}
argumentans 
\edtext{male supponit pag.~139\lbrack:\rbrack\
\textit{patet quod pondus \textit{T} est ad vim qua idem \textit{T} nititur et comprimit planum \textit{IC},}
\edtext{\textit{ut \textit{IC}}}{\lemma{}\Bfootnote{\textit{ut} \bars{ut \textit{streicht Hrsg.}} \textit{IC}~\textit{L}}}
\textit{ad \textit{LC}.}
Et contra 
\indexN{\textls{Stevin} (Stevinus), Simon 1548\textendash1620}Stevinum et
\edtext{\indexN[Herigone]{\textls{Hérigone}, Pierre 1580\textendash1643}Herigonum pag.~141}{%
\lemma{Herigonum}%
\Bfootnote{%
\textit{(1)}~male
\textit{(2)}~pag.~141~\textit{L}}}
male dicit: \textit{tunc pondus \textit{T} dum moveri niteretur per duas rectas inclinatas \textit{CK},
et \textit{CG} cogeretur moveri aut nisum\indexS{nisus} exercere per diagonalem \textit{CO} 
secantem angulum \textit{GCK} bifariam.}}{%
\lemma{male \lbrack...\rbrack\ \textit{bifariam}}%
\Cfootnote{\indexN{\textls{Varignon}, Pierre de 1654\textendash1722}\textsc{Varignon},
\cite{02028}\textit{Examen}, S.\;97 mit Änderungen.
Vgl. \indexN{\textls{Borelli} (Borellus), Giovanni Alfonso 1608\textendash1679}\textsc{Borelli},
\cite{02099}\textit{De motu animalium}, Digressio, S.\;139 und 141, sowie
\indexN[Herigone]{\textls{Hérigone}, Pierre 1580\textendash1643}\textsc{P.\,Hérigone},
\cite{02127}\textit{Cursus mathematicus}, Tome~III, Paris 1634, \textit{Mechanica}, prop.~XII, S.\;313\,f.
und
\indexN{\textls{Stevin} (Stevinus), Simon 1548\textendash1620}\textsc{S.~Stevin}, 
\cite{00196}\textit{De Beghinselen der Weeghconst}, Leiden 1586, 
Buch~1, Teil~2, Theor.~XVIII Prop.~XXVII, S.\;60\,f. (\cite{02084}\textit{SPW} I, S.\;212\textendash215).}}
\edtext{Item male dicit: quod \textit{vis quam patitur planum\indexS{vis quam patitur planum} \textit{CO} 
a compressione ponderis \textit{T} aequalis est viribus ambarum potentiarum\indexS{potentia} \textit{R} et \textit{S}, 
quae sustinendo idem pondus in tali situ, plani \textit{CO} inclinati\indexS{planum inclinatum} vicem supplent.}}{%
\lemma{Item \lbrack...\rbrack\ \textit{supplent}}%
\Cfootnote{\indexN{\textls{Varignon}, Pierre de 1654\textendash1722}\textsc{Varignon},
\cite{02028}\textit{Examen}, S.\;98 mit Änderungen.}}%
\edlabel{37_05_217-218_5b}%
\pend
